\documentclass[11pt]{article}

\usepackage[margin=1in]{geometry}
\usepackage[T1]{fontenc}
\usepackage[utf8]{inputenc}
\usepackage{lmodern}
\usepackage{microtype}
\usepackage{amsmath, amssymb, amsthm}
\usepackage{mathtools}
\usepackage{xcolor}
\usepackage{hyperref}
\usepackage{parskip}

\definecolor{linknavy}{HTML}{1A4FA0}
\hypersetup{
  colorlinks=true,
  urlcolor=linknavy,
  linkcolor=linknavy,
  pdftitle={Notes on Vakil -- The Rising Sea},
  pdfauthor={Vandan Pokal}
}

% Standard theorem environments
\newtheorem{theorem}{Theorem}[section]
\newtheorem{lemma}[theorem]{Lemma}
\newtheorem{proposition}[theorem]{Proposition}
\newtheorem{corollary}[theorem]{Corollary}

\theoremstyle{definition}
\newtheorem{definition}[theorem]{Definition}
\newtheorem{example}[theorem]{Example}
\newtheorem{exercise}[theorem]{Exercise}

\theoremstyle{remark}
\newtheorem{remark}[theorem]{Remark}

% Common AG / category-theory shorthand
\newcommand{\Spec}{\operatorname{Spec}}
\newcommand{\Proj}{\operatorname{Proj}}
\newcommand{\Hom}{\operatorname{Hom}}
\newcommand{\End}{\operatorname{End}}
\newcommand{\Aut}{\operatorname{Aut}}
\newcommand{\im}{\operatorname{im}}
\newcommand{\rk}{\operatorname{rank}}
\newcommand{\id}{\operatorname{id}}
\newcommand{\op}{\mathrm{op}}
\newcommand{\Set}{\mathbf{Set}}
\newcommand{\Ring}{\mathbf{Ring}}
\newcommand{\CRing}{\mathbf{CRing}}
\newcommand{\Top}{\mathbf{Top}}
\newcommand{\Sch}{\mathbf{Sch}}
\newcommand{\Mod}{\mathbf{Mod}}
\renewcommand{\AA}{\mathbb{A}}
\newcommand{\PP}{\mathbb{P}}
\newcommand{\ZZ}{\mathbb{Z}}
\newcommand{\QQ}{\mathbb{Q}}
\newcommand{\RR}{\mathbb{R}}
\newcommand{\CC}{\mathbb{C}}
\newcommand{\Ff}{\mathbb{F}}            % $\Ff$ for finite fields
\newcommand{\OO}{\mathcal{O}}
\newcommand{\Fcal}{\mathcal{F}}         % $\Fcal$ for sheaf F
\newcommand{\Gcal}{\mathcal{G}}
\newcommand{\into}{\hookrightarrow}
\newcommand{\onto}{\twoheadrightarrow}
\newcommand{\iso}{\xrightarrow{\sim}}

\title{Notes on Vakil's \emph{The Rising Sea} \\[2pt]
       \large Foundations of Algebraic Geometry}
\author{Vandan Pokal}
\date{Started May 2026}

\begin{document}
\maketitle

\tableofcontents

\section*{Reading log}

\begin{itemize}
  \item \textbf{May 2026 ---} started; setting up notation and motivation.
\end{itemize}

% =====================================================================
\section{Introduction and motivation}

\textit{What is the book trying to do?} Replace the classical
varieties-over-a-field treatment with the Grothendieck-school
foundation: schemes, sheaves of modules, cohomology. The point is to
build the vocabulary in which modern algebraic geometry is actually
written.

\textit{Vakil's stylistic bet.} Develop the abstract machinery first
and let the geometric intuition arrive via examples that are forced
through the machinery, rather than the other way around. ``The rising
sea'' is Grothendieck's metaphor: do not chisel through the rock,
flood it.

\subsection*{What I want from this pass}

\begin{itemize}
  \item Comfort with the language: schemes, morphisms, quasi-coherent sheaves, cohomology of sheaves.
  \item Worked exercises through at least the schemes / quasi-coherent sheaves chapters.
  \item Enough fluency to read a paper that uses this language without translation.
\end{itemize}

% =====================================================================
\section{Category theory background (Ch.~1)}

% Replace this with actual notes as you go.
\textit{Notes go here.}

\subsection{Categories, functors, natural transformations}

% =====================================================================
\section{Sheaves (Ch.~2)}

\textit{Notes go here.}

% =====================================================================
\section{Toward affine schemes (Ch.~3)}

\textit{Notes go here.}

% =====================================================================
\section{Worked exercises}

Numbered to match Vakil's numbering. Add as worked.

% \begin{exercise}[Vakil 1.x.y]
%   Statement.
% \end{exercise}
% \begin{proof}
%   Proof.
% \end{proof}

% =====================================================================
\section*{Open questions / things to come back to}

\begin{itemize}
  \item ---
\end{itemize}

\end{document}
